\documentclass[12pt,a4paper]{article}

% ---- Encoding & Font ----
\usepackage{iftex}
\ifPDFTeX
  \usepackage[utf8]{inputenc}
  \usepackage[T1]{fontenc}
\else
  \usepackage{fontspec}  % XeTeX/LuaTeX (tectonic): Unicode nativo (·, —, ¿, ª, §)
\fi
\usepackage[spanish]{babel}
\usepackage{csquotes}

% ---- Page layout ----
\usepackage[top=2.5cm, bottom=2.5cm, left=2.5cm, right=2.5cm]{geometry}
\usepackage{setspace}
\onehalfspacing

% ---- Math ----
\usepackage{amsmath, amssymb, amsthm}
\usepackage{mathtools}
\usepackage{physics}
\usepackage{esint}  % \oiint

% ---- Graphics ----
\usepackage{graphicx}
\usepackage{tikz}
\usepackage{float}
\usepackage{caption}

% ---- Tables ----
\usepackage{array}
\usepackage{booktabs}
\usepackage{tabularx}
\usepackage{colortbl}

% ---- Colours ----
\usepackage{xcolor}
\definecolor{notebg}{HTML}{FFF3CD}
\definecolor{noteborder}{HTML}{FFC107}
\definecolor{boxred}{HTML}{F8D7DA}
\definecolor{boxredborder}{HTML}{F5C6CB}

% ---- Boxes ----
\usepackage{mdframed}
\usepackage{tcolorbox}
\tcbuselibrary{most}

\newtcolorbox{keybox}[1][]{
  colback=blue!5,
  colframe=blue!70!black,
  arc=4pt,
  boxrule=1pt,
  fonttitle=\bfseries,
  title={#1}
}

\newtcolorbox{warnbox}[1][]{
  colback=boxred,
  colframe=boxredborder,
  coltitle=black!75,
  arc=4pt,
  boxrule=1pt,
  fonttitle=\bfseries,
  title={#1}
}

\newtcolorbox{notebox}[1][]{
  colback=notebg,
  colframe=noteborder,
  coltitle=black!75,
  arc=4pt,
  boxrule=1pt,
  fonttitle=\bfseries,
  title={#1}
}

% ---- Diagramas (gráficas y esquemas) ----
\usepackage{pgfplots}
\pgfplotsset{compat=1.18}
\usetikzlibrary{babel}  % babel español activa `>`; esto evita romper las flechas `->`
% courses/ElecMag2024/Clases/assets/tikzstyles.tex
% ---------------------------------------------------------------------------
% Estilos compartidos para los diagramas del curso Electromagnetismo (2024).
% Fuente única del "look" visual (colores, grosores, escalas) para que todas
% las figuras (TikZ / pgfplots / CircuiTikz) sean consistentes entre clases.
%
% Es una COPIA de courses/Fisica3-2015/Clases/assets/tikzstyles.tex: mismo
% dominio, misma paleta. Copia y no symlink para que cada curso sea
% autocontenido y la paleta de Física III pueda evolucionar aparte.
%
% Uso: la clase debe cargar antes `tikz` y `pgfplots` (y `circuitikz` si la
%      clase tiene circuitos), y luego % courses/ElecMag2024/Clases/assets/tikzstyles.tex
% ---------------------------------------------------------------------------
% Estilos compartidos para los diagramas del curso Electromagnetismo (2024).
% Fuente única del "look" visual (colores, grosores, escalas) para que todas
% las figuras (TikZ / pgfplots / CircuiTikz) sean consistentes entre clases.
%
% Es una COPIA de courses/Fisica3-2015/Clases/assets/tikzstyles.tex: mismo
% dominio, misma paleta. Copia y no symlink para que cada curso sea
% autocontenido y la paleta de Física III pueda evolucionar aparte.
%
% Uso: la clase debe cargar antes `tikz` y `pgfplots` (y `circuitikz` si la
%      clase tiene circuitos), y luego % courses/ElecMag2024/Clases/assets/tikzstyles.tex
% ---------------------------------------------------------------------------
% Estilos compartidos para los diagramas del curso Electromagnetismo (2024).
% Fuente única del "look" visual (colores, grosores, escalas) para que todas
% las figuras (TikZ / pgfplots / CircuiTikz) sean consistentes entre clases.
%
% Es una COPIA de courses/Fisica3-2015/Clases/assets/tikzstyles.tex: mismo
% dominio, misma paleta. Copia y no symlink para que cada curso sea
% autocontenido y la paleta de Física III pueda evolucionar aparte.
%
% Uso: la clase debe cargar antes `tikz` y `pgfplots` (y `circuitikz` si la
%      clase tiene circuitos), y luego \input{../assets/tikzstyles.tex}
% (compilar desde el directorio de la clase para que la ruta relativa resuelva).
% ---------------------------------------------------------------------------

% ---- Paleta (alineada con las cajas del preámbulo: keybox/notebox/warnbox) ----
\definecolor{figblue}{HTML}{1F4E79}   % azul principal (líneas, keybox)
\definecolor{figred}{HTML}{C0392B}    % rojo de énfasis (warnbox)
\definecolor{figamber}{HTML}{B8860B}  % ámbar (notebox)
\definecolor{figgray}{HTML}{555555}   % auxiliares, ejes, guías

% ---- CircuiTikz: escala y grosor de los componentes ----
% Guarda \ifdefined: la electrostática (Clases 1-14) no carga `circuitikz`, y
% sin la guarda el \input revienta con `Undefined control sequence \ctikzset`.
% Cuando el curso llegue a circuitos, el bloque se activa solo.
\ifdefined\ctikzset
  \ctikzset{
    bipoles/thickness=1,
    resistors/scale=0.9,
    capacitors/scale=0.9,
  }
\fi

% ---- TikZ: estilos reutilizables ----
% Nota: las puntas se declaran como `-{Latex[...]}` y no como `->`. La forma
% con llaves NO contiene un `>` literal, así que es inmune al `>` activo que
% introduce babel español (ver CLAUDE.md §7.3.3).
\usetikzlibrary{arrows.meta,calc,angles,quotes}

\tikzset{
  figline/.style={figblue, line width=0.9pt},
  figaux/.style={figgray, dashed, line width=0.6pt},
  % vectores
  figvec/.style={figblue, line width=0.9pt, -{Latex[length=2.2mm,width=1.6mm]}},
  figvecres/.style={figred, line width=1.3pt, -{Latex[length=2.6mm,width=1.9mm]}},
  figvecaux/.style={figgray, line width=0.7pt, -{Latex[length=1.8mm,width=1.3mm]}},
  figaxis/.style={figgray, line width=0.6pt, -{Latex[length=2mm,width=1.4mm]}},
  % cargas puntuales
  figcarga/.style={circle, inner sep=0pt, minimum size=5.4pt, draw=none},
  figpos/.style={figcarga, fill=figred},
  figneg/.style={figcarga, fill=figblue},
  figneutra/.style={figcarga, fill=figgray},
  % etiquetas
  figlbl/.style={font=\small},
  figlblaux/.style={font=\footnotesize, figgray},
}

% ---- pgfplots: estilo base de las gráficas ----
\pgfplotsset{
  emfig/.style={
    width=11cm, height=6.5cm,
    axis lines=left,
    xlabel style={font=\small},
    ylabel style={font=\small},
    tick label style={font=\footnotesize},
    every axis plot/.append style={line width=1pt},
    grid=major,
    grid style={figgray!20},
    legend cell align=left,
    legend style={font=\footnotesize, draw=figgray!40},
  },
}

% (compilar desde el directorio de la clase para que la ruta relativa resuelva).
% ---------------------------------------------------------------------------

% ---- Paleta (alineada con las cajas del preámbulo: keybox/notebox/warnbox) ----
\definecolor{figblue}{HTML}{1F4E79}   % azul principal (líneas, keybox)
\definecolor{figred}{HTML}{C0392B}    % rojo de énfasis (warnbox)
\definecolor{figamber}{HTML}{B8860B}  % ámbar (notebox)
\definecolor{figgray}{HTML}{555555}   % auxiliares, ejes, guías

% ---- CircuiTikz: escala y grosor de los componentes ----
% Guarda \ifdefined: la electrostática (Clases 1-14) no carga `circuitikz`, y
% sin la guarda el \input revienta con `Undefined control sequence \ctikzset`.
% Cuando el curso llegue a circuitos, el bloque se activa solo.
\ifdefined\ctikzset
  \ctikzset{
    bipoles/thickness=1,
    resistors/scale=0.9,
    capacitors/scale=0.9,
  }
\fi

% ---- TikZ: estilos reutilizables ----
% Nota: las puntas se declaran como `-{Latex[...]}` y no como `->`. La forma
% con llaves NO contiene un `>` literal, así que es inmune al `>` activo que
% introduce babel español (ver CLAUDE.md §7.3.3).
\usetikzlibrary{arrows.meta,calc,angles,quotes}

\tikzset{
  figline/.style={figblue, line width=0.9pt},
  figaux/.style={figgray, dashed, line width=0.6pt},
  % vectores
  figvec/.style={figblue, line width=0.9pt, -{Latex[length=2.2mm,width=1.6mm]}},
  figvecres/.style={figred, line width=1.3pt, -{Latex[length=2.6mm,width=1.9mm]}},
  figvecaux/.style={figgray, line width=0.7pt, -{Latex[length=1.8mm,width=1.3mm]}},
  figaxis/.style={figgray, line width=0.6pt, -{Latex[length=2mm,width=1.4mm]}},
  % cargas puntuales
  figcarga/.style={circle, inner sep=0pt, minimum size=5.4pt, draw=none},
  figpos/.style={figcarga, fill=figred},
  figneg/.style={figcarga, fill=figblue},
  figneutra/.style={figcarga, fill=figgray},
  % etiquetas
  figlbl/.style={font=\small},
  figlblaux/.style={font=\footnotesize, figgray},
}

% ---- pgfplots: estilo base de las gráficas ----
\pgfplotsset{
  emfig/.style={
    width=11cm, height=6.5cm,
    axis lines=left,
    xlabel style={font=\small},
    ylabel style={font=\small},
    tick label style={font=\footnotesize},
    every axis plot/.append style={line width=1pt},
    grid=major,
    grid style={figgray!20},
    legend cell align=left,
    legend style={font=\footnotesize, draw=figgray!40},
  },
}

% (compilar desde el directorio de la clase para que la ruta relativa resuelva).
% ---------------------------------------------------------------------------

% ---- Paleta (alineada con las cajas del preámbulo: keybox/notebox/warnbox) ----
\definecolor{figblue}{HTML}{1F4E79}   % azul principal (líneas, keybox)
\definecolor{figred}{HTML}{C0392B}    % rojo de énfasis (warnbox)
\definecolor{figamber}{HTML}{B8860B}  % ámbar (notebox)
\definecolor{figgray}{HTML}{555555}   % auxiliares, ejes, guías

% ---- CircuiTikz: escala y grosor de los componentes ----
% Guarda \ifdefined: la electrostática (Clases 1-14) no carga `circuitikz`, y
% sin la guarda el \input revienta con `Undefined control sequence \ctikzset`.
% Cuando el curso llegue a circuitos, el bloque se activa solo.
\ifdefined\ctikzset
  \ctikzset{
    bipoles/thickness=1,
    resistors/scale=0.9,
    capacitors/scale=0.9,
  }
\fi

% ---- TikZ: estilos reutilizables ----
% Nota: las puntas se declaran como `-{Latex[...]}` y no como `->`. La forma
% con llaves NO contiene un `>` literal, así que es inmune al `>` activo que
% introduce babel español (ver CLAUDE.md §7.3.3).
\usetikzlibrary{arrows.meta,calc,angles,quotes}

\tikzset{
  figline/.style={figblue, line width=0.9pt},
  figaux/.style={figgray, dashed, line width=0.6pt},
  % vectores
  figvec/.style={figblue, line width=0.9pt, -{Latex[length=2.2mm,width=1.6mm]}},
  figvecres/.style={figred, line width=1.3pt, -{Latex[length=2.6mm,width=1.9mm]}},
  figvecaux/.style={figgray, line width=0.7pt, -{Latex[length=1.8mm,width=1.3mm]}},
  figaxis/.style={figgray, line width=0.6pt, -{Latex[length=2mm,width=1.4mm]}},
  % cargas puntuales
  figcarga/.style={circle, inner sep=0pt, minimum size=5.4pt, draw=none},
  figpos/.style={figcarga, fill=figred},
  figneg/.style={figcarga, fill=figblue},
  figneutra/.style={figcarga, fill=figgray},
  % etiquetas
  figlbl/.style={font=\small},
  figlblaux/.style={font=\footnotesize, figgray},
}

% ---- pgfplots: estilo base de las gráficas ----
\pgfplotsset{
  emfig/.style={
    width=11cm, height=6.5cm,
    axis lines=left,
    xlabel style={font=\small},
    ylabel style={font=\small},
    tick label style={font=\footnotesize},
    every axis plot/.append style={line width=1pt},
    grid=major,
    grid style={figgray!20},
    legend cell align=left,
    legend style={font=\footnotesize, draw=figgray!40},
  },
}


% ---- Hyperlinks & TOC ----
\usepackage[hidelinks]{hyperref}
\usepackage{bookmark}

% ---- Enumerate ----
\usepackage{enumitem}

% ---- Miscellaneous ----
\usepackage{icomma}
\usepackage{siunitx}
\usepackage{textcomp}
\usepackage[bottom]{footmisc}

% ---- Title ----
\title{\textbf{Clase 5: Ecuaciones de Poisson y Laplace} \\[0.3em]
        \large{Cierre del desarrollo multipolar y teorema de unicidad}}
\author{Transcripto y expandido --- Curso Electromagnetismo (OpenFING)}
\date{}

\begin{document}

\maketitle
\thispagestyle{empty}
\newpage

\tableofcontents
\newpage


\section{El desarrollo multipolar, término a término}

La clase retoma exactamente donde quedó la anterior. Para una distribución
continua localizada $\rho(\vec r\,')$ en un volumen $B$, con el origen
\textbf{dentro} de $B$ y $r \gg a$ (siendo $a$ el tamaño típico), el potencial
exacto es

\[
\phi(\vec r) = \frac{1}{4\pi\varepsilon_0}\int_B \frac{\rho(\vec r\,')}{|\vec r - \vec r\,'|}\,dV'
\]

y el desarrollo de Taylor obtenido en la Clase 4 es

\[
\frac{1}{|\vec r - \vec r\,'|} = \frac{1}{r} + \frac{\vec r\cdot\vec r\,'}{r^{3}}
+ \frac{1}{2}\left[\frac{3(\vec r\cdot\vec r\,')^{2}}{r^{5}} - \frac{r'^{2}}{r^{3}}\right]
+ \mathcal{O}(r'^{3})
\]

\begin{notebox}[Dónde cortar]
El desarrollo se puede llevar al orden que se quiera; dónde cortar lo fija la
precisión buscada, como en cualquier Taylor. Acá se corta en $r'^2$, que es lo
mismo que decir en $1/r^3$: el primer término es de orden $1/r$, el segundo de
orden $1/r^2$ (hay un $r$ arriba y un $r^3$ abajo) y el tercero de orden $1/r^3$.
\end{notebox}

\subsection{El problema de sacar $\vec r$ fuera de la integral}

La estrategia es meter el desarrollo en la integral y \textbf{sacar afuera todo
lo que dependa de $\vec r$}, porque $\vec r$ no se integra. Con los dos primeros
términos sale solo:

\[
\phi(\vec r) = \frac{1}{4\pi\varepsilon_0}\frac{1}{r}\int_B \rho(\vec r\,')\,dV'
\;+\; \frac{1}{4\pi\varepsilon_0}\frac{\vec r}{r^{3}}\cdot\int_B \rho(\vec r\,')\,\vec r\,'\,dV'
\;+\;\cdots
\]

Pero el tercer término tiene un producto escalar \textbf{al cuadrado},
$(\vec r\cdot\vec r\,')^2$, y ahí la factorización deja de ser obvia: no hay
manera de sacar $\vec r$ como un factor vectorial limpio.

\subsection{La delta de Kronecker como herramienta}

La salida es pasar a componentes. Se introduce la \textbf{delta de Kronecker}:

\[
\delta_{ij} = \begin{cases} 1 & \text{si } i = j\\ 0 & \text{si } i \neq j\end{cases}
\]

Con ella, el producto escalar y su cuadrado se escriben por componentes:

\[
\vec r\cdot\vec r\,' = \sum_{i=1}^{3} r_i r'_i
\qquad\Longrightarrow\qquad
(\vec r\cdot\vec r\,')^2 = \sum_{i,j=1}^{3} r_i r'_i\, r_j r'_j
\]

Y ---éste es el truco--- el término $r'^2/r^3$, que a primera vista no tiene la
misma forma, \textbf{se puede forzar a tenerla}. Notando que

\[
r'^2 = \sum_{i,j=1}^{3} r'_i r'_j\,\delta_{ij} = \sum_{i=1}^{3} r_i'^2
\qquad\text{y análogamente}\qquad
r^2 = \sum_{i,j=1}^{3} r_i r_j\,\delta_{ij}
\]

se puede escribir
$\dfrac{r'^2}{r^3} = \dfrac{r'^2\,r^2}{r^5} = \sum_{i,j}\dfrac{r_i r_j}{r^5}\,r'^2\,\delta_{ij}$,
que ya tiene el mismo factor $r_i r_j / r^5$ que el otro término.

\begin{notebox}[Por qué complicar lo que parecía simple]
Escribir $r'^2$ como una suma doble con una delta es más largo, pero es lo único
que permite \textbf{sacar toda la dependencia en $\vec r$ fuera de la integral en
un solo factor común}. Ése es el objetivo del ejercicio.
\end{notebox}

Juntando los dos pedazos bajo la misma suma:

\[
\text{3.\textsuperscript{er} término} =
\frac{1}{4\pi\varepsilon_0}\cdot\frac{1}{2}\sum_{i,j=1}^{3}\frac{r_i r_j}{r^{5}}
\int_B \rho(\vec r\,')\left[3 r'_i r'_j - r'^2\,\delta_{ij}\right]dV'
\]

\subsection{Los tres momentos y el potencial}

Cada integral que quedó es una propiedad \textbf{de la distribución sola}, sin
ninguna referencia al punto de observación. Se les da nombre:

\begin{table}[H]
  \centering
  \begin{tabularx}{\textwidth}{@{}l >{$\displaystyle}X<{$} c@{}}
    \toprule
    \textbf{Momento} & \multicolumn{1}{c}{\textbf{Definición}} & \textbf{N.\textsuperscript{o} de números} \\
    \midrule
    Carga total & Q = \int_B \rho(\vec r\,')\,dV' & 1 \\[0.5em]
    Momento dipolar & \vec p = \int_B \rho(\vec r\,')\,\vec r\,'\,dV' & 3 \\[0.5em]
    Momento cuadrupolar & Q_{ij} = \int_B \rho(\vec r\,')\left[3 r'_i r'_j - r'^2\delta_{ij}\right]dV' & 6 \\
    \bottomrule
  \end{tabularx}
\end{table}

\begin{notebox}[El dipolar es una generalización, no algo nuevo]
Es el mismo $\vec p = Q\vec L$ de la clase anterior: allá se definió para dos
cargas puntuales, acá para una distribución continua, pero es la misma idea.
\end{notebox}

Con esto el potencial queda escrito como una serie de términos con potencias
crecientes de $1/r$:

\[
\boxed{\;\phi(\vec r) = \frac{1}{4\pi\varepsilon_0}\frac{Q}{r}
\;+\; \frac{1}{4\pi\varepsilon_0}\frac{\vec p\cdot\vec r}{r^{3}}
\;+\; \frac{1}{2}\sum_{i,j=1}^{3}\frac{1}{4\pi\varepsilon_0}\frac{r_i r_j\,Q_{ij}}{r^{5}}
\;+\; \mathcal{O}\!\left(\frac{1}{r^{4}}\right)\;}
\]

\begin{warnbox}[El tercer término es de orden $1/r^3$]
Aunque el denominador diga $r^5$: los dos $r_i r_j$ del numerador se comen dos
potencias. Es la misma contabilidad que en el segundo término, donde el $r^3$ de
abajo y el $\vec r$ de arriba dejan $1/r^2$. La regla general: cada $\vec r\,'$
extra en el numerador trae un $1/r$ extra en el denominador --- que es justamente
por qué da lo mismo pensarlo como un desarrollo en $r'$ chico o en $r$ grande: lo
que se desarrolla es el \textbf{cociente}.
\end{warnbox}

El desarrollo sigue: el término siguiente se llama \textbf{momento octopolar} y
tiene tres índices, el que le sigue cuatro, y así. El docente corta en $1/r^3$
«porque después las cuentas empiezan a ser demoníacas», aclarando que
conceptualmente se podría llegar tan lejos como se quiera.

\section{Qué dice el desarrollo}

\subsection{La jerarquía de comportamientos}

El desarrollo formaliza lo que se había anticipado cualitativamente en la Clase 4.

\begin{figure}[H]
  \centering
  % Clase 5 — la jerarquia que el desarrollo formaliza: que domina a gran distancia
% segun cuales momentos se anulen. Configuraciones minimas y su decaimiento.
\begin{tikzpicture}[scale=1]

% ---- monopolo ----
\begin{scope}
  \node[figpos] at (0,0) {};
  \node[figlbl] at (0,-1.3) {monopolo};
  \node[figlbl] at (0,-1.9) {$Q\neq0$};
  \node[figlbl, figblue] at (0,-2.5) {$\phi\sim 1/r$};
\end{scope}

% ---- dipolo ----
\begin{scope}[xshift=4.4cm]
  \node[figpos] at (-0.45,0) {};
  \node[figneg] at (0.45,0) {};
  \node[figlblaux] at (-0.45,0.42) {$+$};
  \node[figlblaux] at (0.45,0.42) {$-$};
  \node[figlbl] at (0,-1.3) {dipolo};
  \node[figlbl] at (0,-1.9) {$Q=0$, $\vec p\neq0$};
  \node[figlbl, figblue] at (0,-2.5) {$\phi\sim 1/r^{2}$};
\end{scope}

% ---- cuadrupolo ----
\begin{scope}[xshift=9.6cm]
  \node[figpos] at (-0.45,0.45) {};
  \node[figneg] at (0.45,0.45) {};
  \node[figneg] at (-0.45,-0.45) {};
  \node[figpos] at (0.45,-0.45) {};
  \node[figlblaux] at (-0.82,0.45) {$+$};
  \node[figlblaux] at (0.82,0.45) {$-$};
  \node[figlblaux] at (-0.82,-0.45) {$-$};
  \node[figlblaux] at (0.82,-0.45) {$+$};
  \node[figlbl] at (0,-1.3) {cuadrupolo};
  \node[figlbl] at (0,-1.9) {$Q=0$, $\vec p=0$};
  \node[figlbl, figblue] at (0,-2.5) {$\phi\sim 1/r^{3}$};
\end{scope}

% flecha de jerarquia
\draw[figvecaux] (1.3,-0.62) -- (3.1,-0.62);
\draw[figvecaux] (5.7,-0.62) -- (8.1,-0.62);
\node[figlblaux, fill=white, inner sep=1.5pt] at (2.2,-0.34) {si $Q=0$};
\node[figlblaux, fill=white, inner sep=1.5pt] at (6.9,-0.34) {si adem\'as $\vec p=0$};

\node[figlbl, align=center] at (4.8,-3.35)
  {Cada orden entra reci\'en cuando los anteriores se anulan;\\[0.2em]
   la serie sigue con el momento octopolar ($\phi\sim 1/r^{4}$) y as\'i.};

\end{tikzpicture}

  \caption*{\small\itshape Cada término manda sólo cuando los anteriores se
  anulan; la configuración mínima que lo realiza va perdiendo carga neta y
  después momento dipolar.}
\end{figure}

\begin{table}[H]
  \centering
  \begin{tabular}{@{}ll@{}}
    \toprule
    \textbf{Si\ldots} & \textbf{el comportamiento dominante a gran distancia es\ldots} \\
    \midrule
    $Q \neq 0$ & el de una \textbf{carga puntual} ($\sim 1/r$) \\
    $Q = 0$ & el de un \textbf{dipolo} ($\sim 1/r^2$) \\
    $Q = 0$ y $\vec p = 0$ & el de un \textbf{cuadrupolo} ($\sim 1/r^3$) \\
    \bottomrule
  \end{tabular}
\end{table}

\subsection{Dependencia del origen de coordenadas}

Hay algo incómodo en estas expresiones: \textbf{dependen del sistema de
coordenadas}, porque todas hacen referencia al punto $O$. Si se cambia el origen,
$\vec p$ cambia.

\begin{notebox}[Eso no es un defecto, es la naturaleza del método]
Es un desarrollo de Taylor, y un Taylor se hace alrededor de un punto que uno
elige. \textbf{La suma completa de la serie no cambia}; lo que cambia es cómo se
reparte entre los términos.

Sí hay una restricción real: \textbf{$O$ tiene que estar dentro de $B$}, para que
$r'$ sea del orden del tamaño del sistema. Con el origen lejos, $r'$ ya no es
chico y el desarrollo no es en el parámetro que se quería.
\end{notebox}

Pero hay una propiedad importante: \textbf{los momentos de un dado orden se
vuelven independientes del origen cuando todos los de orden menor son nulos.} El
caso concreto que se demuestra es el del momento dipolar.

\begin{figure}[H]
  \centering
  % Clase 5 — los momentos dependen del origen elegido. Geometria del dibujo con
% que el docente prueba que p es independiente del origen cuando Q=0.
\begin{tikzpicture}[scale=1.2]

  % volumen B
  \fill[figblue!12] plot[smooth cycle, tension=0.85]
    coordinates {(-1.05,0.5) (-0.1,1.05) (0.95,0.6) (1.05,-0.4) (0.1,-0.95) (-1.0,-0.45)};
  \draw[figline] plot[smooth cycle, tension=0.85]
    coordinates {(-1.05,0.5) (-0.1,1.05) (0.95,0.6) (1.05,-0.4) (0.1,-0.95) (-1.0,-0.45)};
  \node[figlbl] at (-1.5,1.05) {$B$};

  % punto fuente
  \node[figneutra] (s) at (0.5,0.35) {};
  \node[figlblaux, fill=white, inner sep=1.5pt] at (1.62,0.72) {$\rho(\vec r\,')$};
  \draw[figvecaux] (1.3,0.62) -- (0.62,0.44);

  % origen O
  \node[figneutra] (O) at (-0.7,-0.2) {};
  \node[figlbl, fill=figblue!12, inner sep=1.5pt] at (-1.0,-0.02) {$O$};
  \draw[figvec] (O) -- (s);
  \node[figlbl, fill=figblue!12, inner sep=1.5pt] at (-0.12,-0.08) {$\vec r\,'$};

  % origen Obarra
  \node[figneutra] (Ob) at (0.55,-1.55) {};
  \node[figlbl] at (0.55,-1.88) {$\bar O$};
  \draw[figvecres] (Ob) -- (s);
  \node[figlbl, figred, fill=white, inner sep=1.5pt] at (0.92,-0.62) {$\vec r\,''$};

  % vector entre origenes
  \draw[figvecaux] (O) -- (Ob);
  \node[figlblaux, fill=white, inner sep=1.5pt] at (-0.72,-1.02) {$\vec{O\bar O}$};

  \node[figlbl, align=center] at (0.0,-2.85)
    {$\vec r\,''=\vec r\,'-\vec{O\bar O}
      \ \Longrightarrow\ \vec p_{\bar O}=\vec p_{O}-\vec{O\bar O}\,Q$\\[0.25em]
     Con $Q=0$ los dos or\'igenes dan el mismo $\vec p$.};

\end{tikzpicture}

  \caption*{\small\itshape El mismo punto fuente medido desde dos orígenes. La
  diferencia entre ambos momentos dipolares es exactamente $\vec{O\bar O}\,Q$, y
  por eso desaparece cuando la carga neta es nula.}
\end{figure}

\textbf{Demostración (dipolar).} Sean dos orígenes, $O$ y $\bar O$, y llámese
$\vec{O\bar O}$ al vector que va de uno a otro. Un mismo punto fuente se mide como
$\vec r\,'$ desde $O$ y como $\vec r\,''$ desde $\bar O$, con
$\vec r\,'' = \vec r\,' - \vec{O\bar O}$. El elemento de volumen no cambia (es un
mero cambio de origen), así que

\[
\vec p_{\bar O} = \int_B \rho(\vec r\,')\,\vec r\,''\,dV'
= \int_B \rho(\vec r\,')\,\vec r\,'\,dV' \;-\; \vec{O\bar O}\int_B \rho(\vec r\,')\,dV'
= \vec p_{O} - \vec{O\bar O}\,Q
\]

\[
\boxed{\;Q = 0 \quad\Longrightarrow\quad \vec p_{\bar O} = \vec p_{O}\;}
\]

\begin{notebox}[Por qué tenía que ser así]
Si la carga neta es cero, a gran distancia el sistema \emph{se comporta} como un
dipolo, y ese comportamiento observable no puede depender del sistema de
coordenadas que eligió quien hace la cuenta. La fórmula además dice exactamente
cuánto vale la ambigüedad cuando $Q\neq0$: es $\vec{O\bar O}\,Q$.
\end{notebox}

\subsection{Para qué sirve, y la versión discreta}

El argumento es de ingeniería, y el docente lo plantea así: una distribución de
carga complicada, vista \textbf{de lejos}, no necesita guardarse entera. Basta
con unos pocos números:

\begin{itemize}[nosep]
  \item la carga total: \textbf{1} número;
  \item el momento dipolar: \textbf{3} números;
  \item el momento cuadrupolar: \textbf{6} números ---no nueve, porque $Q_{ij}$
        es simétrica: la diagonal más los de fuera de la diagonal.
\end{itemize}

\begin{notebox}[Moraleja de modelado]
Esto sólo sirve si se va a mirar de lejos. De cerca «no hay tutía»: hay que
quedarse con toda la distribución. La idea es guardar la mínima información
compatible con la precisión con la que se va a trabajar.

Ejemplos que se mencionan: una \textbf{antena} vista de lejos se reemplaza por su
momento dipolar (todo esto es estático, pero el mismo desarrollo se hace para
sistemas que dependen del tiempo). En \textbf{ondas gravitacionales}, en cambio,
no hay efecto dipolar y lo que domina es el cuadrupolar, lo que obliga a un
tratamiento más complicado.
\end{notebox}

Todo lo anterior se hizo para una distribución continua, pero funciona igual para
un conjunto discreto de $N$ cargas puntuales: donde había integrales hay sumas, y
donde había densidades hay cargas.

\[
Q = \sum_{k=1}^{N} q_k
\qquad
\vec p = \sum_{k=1}^{N} q_k\,\vec r\,'_k
\qquad
Q_{ij} = \sum_{k=1}^{N} q_k\left[3\,r'_{k,i}\,r'_{k,j} - r_k'^2\,\delta_{ij}\right]
\]

\begin{warnbox}[Aclaración pedida en clase]
Cualquier distribución, para ser descrita \emph{exactamente}, necesita
\textbf{toda} la serie. Un dipolo real ---dos cargas a distancia finita, no
infinitesimal--- tiene también momento cuadrupolar, octopolar, etcétera.
Quedarse en el término dipolar es una decisión de precisión, no una propiedad del
objeto.
\end{warnbox}

\begin{notebox}[Queda pendiente]
La última sección del capítulo 1, la \textbf{delta de Dirac}, se deja
explícitamente «en el tintero» para retomarla más adelante, cuando haga falta.
\end{notebox}

\section{Capítulo 2: por qué hace falta un método nuevo}

Un estudiante podría objetar que la electrostática ya está resuelta, y en parte
tendría razón: \textbf{si las cargas están dadas, con posiciones prefijadas, el
problema está resuelto}. El campo deriva de un potencial y el potencial tiene una
expresión explícita,

\[
\phi(\vec r) = \frac{1}{4\pi\varepsilon_0}\int \frac{\rho(\vec r\,')}{|\vec r - \vec r\,'|}\,dV'
\]

y con $\vec E = -\grad\phi$ se termina. Si la integral es horrible, se discretiza
y se resuelve numéricamente (Simpson, cuadraturas gaussianas, lo que sea) con las
cifras significativas que se quieran. No hay más nada que hacer.

\textbf{Pero hay otro tipo de problema electrostático: aquel en que la posición
de la carga es parte de la incógnita.} El ejemplo de juguete es una \textbf{esfera
conductora colocada en un campo eléctrico externo uniforme}: la esfera se
polariza, la carga se redistribuye en su superficie, y el campo resultante toma
una forma que no se conoce de antemano. La integral de Coulomb no sirve, porque
\textbf{no se sabe dónde están las cargas}.

\begin{figure}[H]
  \centering
  % Clase 5 — el problema motivador del capitulo 2: esfera conductora en un campo
% externo uniforme. La carga se redistribuye sola, asi que no se sabe donde esta
% y la integral de Coulomb no sirve.
\begin{tikzpicture}[scale=1]

% ---------------- panel (a): antes ----------------
\begin{scope}
  \foreach \yy in {-1.5,-0.75,0,0.75,1.5}{
    \draw[figvecaux] (-1.9,\yy) -- (1.9,\yy);
  }
  \node[figlblaux] at (0,1.95) {$\vec E_0$ uniforme};
  \node[figlbl] at (0,-2.35) {(a) campo externo dado};
\end{scope}

% ---------------- panel (b): despues ----------------
\begin{scope}[xshift=6.4cm]
  % lineas de campo que rodean la esfera
  \draw[figvecaux] (-2.3,1.5) -- (2.3,1.5);
  \draw[figvecaux] (-2.3,-1.5) -- (2.3,-1.5);
  \draw[figvecaux] plot[smooth, tension=0.8]
    coordinates {(-2.3,0.75) (-1.2,0.78) (0,1.12) (1.2,0.78) (2.3,0.75)};
  \draw[figvecaux] plot[smooth, tension=0.8]
    coordinates {(-2.3,-0.75) (-1.2,-0.78) (0,-1.12) (1.2,-0.78) (2.3,-0.75)};
  % tramos que entran y salen de la esfera (el interior esta apantallado)
  \draw[figvecaux] (-2.3,0) -- (-0.95,0);
  \draw[figvecaux] (0.95,0) -- (2.3,0);
  % esfera conductora
  \fill[figgray!25] (0,0) circle (0.92);
  \draw[figline] (0,0) circle (0.92);
  \node[figlblaux] at (0,0) {$\vec E=0$};
  % carga inducida: negativa a la izquierda, positiva a la derecha
  \foreach \ang in {150,165,180,195,210}{
    \node[figneg] at (\ang:0.92) {};
  }
  \foreach \ang in {-30,-15,0,15,30}{
    \node[figpos] at (\ang:0.92) {};
  }
  \node[figlbl, figblue, fill=white, inner sep=1.5pt] at (-1.55,0.5) {$\sigma<0$};
  \node[figlbl, figred, fill=white, inner sep=1.5pt] at (1.55,0.5) {$\sigma>0$};
  \node[figlbl] at (0,-2.35) {(b) esfera conductora: se polariza};
\end{scope}

\node[figlbl, align=center] at (3.2,-3.4)
  {La densidad inducida $\sigma$ \textbf{es parte de la inc\'ognita}: no se puede\\[0.2em]
   escribir la integral de Coulomb porque no se sabe d\'onde est\'an las cargas.};

\end{tikzpicture}

  \caption*{\small\itshape El problema que motiva todo el capítulo: entre (a) y
  (b) no cambia el campo aplicado sino quién decide dónde se ubica la carga.}
\end{figure}

\begin{keybox}[Objetivo del capítulo]
Hallar métodos que resuelvan \textbf{a la vez} el campo y la posición de las
cargas, al menos en situaciones relativamente simples. El docente aclara el
alcance: éste es un curso de nivel \textbf{intermedio}, así que no se verá el caso
más general, pero sí se llegará a situaciones relativamente realistas.
\end{keybox}

\begin{notebox}[Alcance: sin dieléctricos]
En este capítulo no hay materiales aislantes que se polaricen localmente. Puede
haber polarización de un metal por circulación de carga, pero no polarización
local del material.
\end{notebox}

\section{La ecuación de Poisson}

\subsection{Deducción}

Las dos ecuaciones que caracterizan la electrostática sin dieléctricos, ya vistas
en el capítulo anterior, son

\[
\text{(1)}\quad \curl \vec E = 0
\qquad\qquad
\text{(2)}\quad \div \vec E = \frac{\rho}{\varepsilon_0}
\]

La primera dice que el campo es \textbf{conservativo}; la segunda es la ley de
Gauss en forma diferencial. Resolver ambas es resolver el problema electrostático.

La (1) se puede aprovechar para simplificar: un campo irrotacional deriva de un
potencial,

\[
\vec E = -\grad\phi
\]

\begin{notebox}[La ganancia de usar el potencial]
Reduce un problema \textbf{vectorial} (tres funciones incógnita) a uno
\textbf{escalar} (una sola).
\end{notebox}

Sustituyendo en (2):

\[
\boxed{\;\div\grad\phi = -\frac{\rho}{\varepsilon_0}\;}
\]

que es la \textbf{ecuación de Poisson}.

\begin{warnbox}[Poisson no es literalmente «el problema general»]
Vale para distribuciones \textbf{continuas} y en los puntos suaves. Donde hay
singularidades ---una carga puntual, o un borde con densidad \textbf{superficial}
de carga, como la que se acumula en la superficie de un conductor--- la ecuación
no vale. El problema completo es entonces: Poisson en el volumen \textbf{más} el
tratamiento de los bordes.
\end{warnbox}

\subsection{El operador laplaciano}

Conviene pensar la composición «gradiente y después divergencia» como \textbf{un
solo operador}, el producto escalar del gradiente consigo mismo. Se define el
\textbf{laplaciano}:

\[
\laplacian\phi \equiv \div\grad\phi
\]

La definición vale para cualquier función escalar y en cualquier sistema de
coordenadas. En cartesianas, partiendo de

\[
\grad\phi = \pdv{\phi}{x}\hat\imath + \pdv{\phi}{y}\hat\jmath + \pdv{\phi}{z}\hat k
\qquad
\div\vec V = \pdv{V_x}{x} + \pdv{V_y}{y} + \pdv{V_z}{z}
\]

y sustituyendo una en la otra ($V_x \to \partial\phi/\partial x$, etc.):

\[
\boxed{\;\laplacian\phi = \pdv[2]{\phi}{x} + \pdv[2]{\phi}{y} + \pdv[2]{\phi}{z}\;}
\]

de modo que Poisson en cartesianas es

\[
\pdv[2]{\phi}{x} + \pdv[2]{\phi}{y} + \pdv[2]{\phi}{z} = -\frac{\rho}{\varepsilon_0}
\]

\subsection{Qué tipo de ecuación es}

Vale la pena clasificarla, porque de eso dependen los métodos:

\begin{itemize}[leftmargin=*]
  \item Tiene derivadas: es una \textbf{ecuación diferencial}.
  \item Tiene derivadas respecto de \textbf{varias} variables (las tres
        coordenadas), no de una sola: es una \textbf{ecuación en derivadas
        parciales}, no una ecuación diferencial \emph{ordinaria}. Ésta es la
        novedad respecto de todo lo que se vio antes.
  \item El laplaciano es un operador \textbf{lineal}: la dependencia en $\phi$ es
        lineal.
  \item Es de \textbf{segundo orden}, y tiene \textbf{segundo miembro}
        ($-\rho/\varepsilon_0$).
\end{itemize}

\begin{notebox}[La analogía con el resorte]
La ecuación del resorte también es lineal, de segundo orden y con segundo
miembro. La única diferencia estructural es que aquélla es ordinaria (derivadas
respecto del tiempo) y ésta es en derivadas parciales (respecto de $x$, $y$, $z$).
\end{notebox}

\section{La ecuación de Laplace}

Cuando en la región de interés \textbf{no hay carga volumétrica} ($\rho = 0$),
Poisson se reduce a la \textbf{ecuación de Laplace}:

\[
\boxed{\;\div\grad\phi = \laplacian\phi = 0\;}
\qquad\Longleftrightarrow\qquad
\pdv[2]{\phi}{x} + \pdv[2]{\phi}{y} + \pdv[2]{\phi}{z} = 0
\]

\begin{notebox}[No es un caso exótico]
Es exactamente el del ejemplo motivador. En la esfera conductora dentro de un
campo externo no hay carga volumétrica \emph{en ningún lado} ---dentro del
conductor no hay (propiedad de la Clase 4) y afuera está el vacío---; toda la
carga vive en la superficie, es decir, en el borde.
\end{notebox}

Laplace es la \textbf{ecuación homogénea asociada} a Poisson: es lo que queda al
«apagar» el segundo miembro. De ahí, por la teoría general de ecuaciones lineales,

\begin{center}
solución general de Poisson $=$ solución general de Laplace $+$ una solución
particular de Poisson
\end{center}

\subsection{Primera propiedad: linealidad}

Si $\phi_1$ y $\phi_2$ son soluciones de Laplace, cualquier combinación lineal
$c_1\phi_1 + c_2\phi_2$ también lo es. La verificación es directa, apoyada en que
gradiente y divergencia son lineales:

\[
\laplacian(c_1\phi_1 + c_2\phi_2) = \div\grad(c_1\phi_1 + c_2\phi_2)
= c_1\laplacian\phi_1 + c_2\laplacian\phi_2 = 0
\]

\begin{keybox}[Consecuencia]
A partir de dos soluciones se construye un conjunto infinito de soluciones. El
\textbf{conjunto de soluciones de Laplace forma un espacio vectorial}.
\end{keybox}

\begin{warnbox}[Cuidado con lo que ese espacio vectorial es y lo que no es]
Es el conjunto de soluciones \textbf{sin imponer condiciones de borde}. Las
distintas soluciones \textbf{pueden no verificar las mismas condiciones de
borde}, y de hecho en general no lo hacen. Resolver una ecuación diferencial
requiere dar la ecuación \textbf{y} las condiciones; acá no son condiciones
\emph{iniciales} (eso remite a un problema temporal) sino \textbf{condiciones de
borde}: cuánto vale el potencial, o su derivada normal, en el borde de la región.
\end{warnbox}

\subsection{Segunda propiedad: unicidad}

\textbf{Planteo geométrico.} Se considera una región $B$ limitada por una
superficie exterior $S_0$ y un conjunto de superficies interiores
$S_1, S_2, \dots$; el problema es resolver Laplace en la región \textbf{entre}
ellas. Físicamente: una cavidad metálica con varios conductores adentro.

\begin{figure}[H]
  \centering
  % Clase 5 — geometria del teorema de unicidad: resolver Laplace en la region B
% (rayada) entre una superficie exterior S0 y superficies interiores S1, S2, ...
% En cada una se da Dirichlet o Von Neumann, y pueden ser de distinto tipo.
\begin{tikzpicture}[scale=1.15]

  % region B, rayada, entre S0 y las Si
  \begin{scope}
    \clip plot[smooth cycle, tension=0.8]
      coordinates {(-2.5,0.9) (-0.6,2.25) (1.9,1.6) (2.5,-0.5) (0.7,-2.1) (-1.9,-1.3)};
    \fill[figgray!12] (-3,-3) rectangle (3,3);
    \foreach \xx in {-3.4,-3.1,...,3.4}{
      \draw[figgray!35, line width=0.35pt] (\xx,-3) -- (\xx+2.6,3);
    }
    % huecos: las regiones interiores no forman parte de B
    \fill[white] (-0.95,0.45) circle (0.68);
    \fill[white] (1.05,-0.75) circle (0.55);
    \fill[white] (0.85,1.05) circle (0.42);
  \end{scope}

  % superficie exterior S0
  \draw[figline, line width=1.1pt] plot[smooth cycle, tension=0.8]
    coordinates {(-2.5,0.9) (-0.6,2.25) (1.9,1.6) (2.5,-0.5) (0.7,-2.1) (-1.9,-1.3)};
  \node[figlbl, figblue] at (-2.35,1.72) {$S_0$};

  % superficies interiores
  \draw[figred, line width=1pt] (-0.95,0.45) circle (0.68);
  \draw[figred, line width=1pt] (1.05,-0.75) circle (0.55);
  \draw[figred, line width=1pt] (0.85,1.05) circle (0.42);
  \node[figlbl, figred] at (-0.95,0.45) {$S_1$};
  \node[figlbl, figred] at (1.05,-0.75) {$S_2$};
  \node[figlbl, figred] at (0.85,1.05) {$S_3$};

  % rotulo de la region
  \node[figlbl, fill=white, inner sep=2pt] at (-0.85,-1.15) {$B$: $\nabla^2\phi=0$};

  % condiciones de borde: una por superficie, pueden ser de distinto tipo
  \node[figlblaux, fill=white, inner sep=1.5pt] at (-0.4,-2.75)
    {Dirichlet: $\phi\big|_{S_i}$ dado \qquad Von Neumann: $\nabla\phi\cdot\hat n\,\big|_{S_i}$ dado};
  \node[figlblaux, fill=white, inner sep=1.5pt, align=left] at (1.15,-2.75)
    {Von Neumann: $\nabla\phi\cdot\hat n\,\big|_{S_i}$ dado};

  \node[figlbl, align=center] at (-0.4,-3.55)
    {En cada superficie se da \emph{una} de las dos, y no tienen\\[0.2em]
     por qu\'e ser del mismo tipo en todas.};

\end{tikzpicture}

  \caption*{\small\itshape La región donde se resuelve Laplace es la rayada: ni
  el exterior de $S_0$ ni el interior de las $S_i$ forman parte de ella.}
\end{figure}

\textbf{Dos tipos de condición de borde}, una por superficie (pueden ser de
distinto tipo en distintas superficies):

\begin{table}[H]
  \centering
  \begin{tabularx}{\textwidth}{@{}c X l@{}}
    \toprule
    \textbf{Tipo} & \textbf{Se da como dato} & \textbf{Nombre} \\
    \midrule
    1 & $\phi\big|_{S_i}$ (no tiene por qué ser constante; en un metal sí lo es,
        por ser equipotencial) & \textbf{Dirichlet} \\[0.4em]
    2 & $\grad\phi\cdot\hat n\,\big|_{S_i}$, es decir la componente normal del
        campo & \textbf{Von Neumann} \\
    \bottomrule
  \end{tabularx}
\end{table}

\begin{notebox}[Alcance]
Se pueden estudiar condiciones más complicadas, o mezclas dentro de una misma
superficie, pero ahí el problema se vuelve más difícil. Acá se consideran estas
dos.
\end{notebox}

\begin{keybox}[Teorema de unicidad]
La ecuación de Laplace en $B$ con esas condiciones de borde tiene solución única,
a menos de una constante aditiva.
\end{keybox}

Si alguna de las superficies tiene condición de \textbf{Dirichlet}, ni siquiera
queda esa constante: el dato la fija. Con condiciones de Von Neumann puras sólo se
dan derivadas, y sumar una constante sigue dando una solución. Esa constante,
además, no es física.

\begin{notebox}[«El teorema de todos los golpes están permitidos»]
Así lo bautiza el docente, porque es la forma útil de pensar cualquier teorema de
unicidad: si hay unicidad y uno encuentra \textbf{una} solución que cumple todas
las condiciones de borde ---por el mecanismo que sea, adivinando, con un truco,
«pegando por debajo de la cintura»--- entonces ganó, porque es \emph{la} solución.

La analogía es el sudoku: si el sudoku está bien diseñado y tiene solución única,
encontrar una solución resuelve el problema; si está mal diseñado y admite varias,
encontrar una no permite afirmar nada. Esto no es una guitarreada: va a ser la
justificación de los métodos de las clases siguientes.
\end{notebox}

\textbf{Demostración (por el absurdo).} Supóngase que hay dos soluciones $\phi_1$
y $\phi_2$ del problema \textbf{completo} (Laplace en $B$ más las condiciones de
borde en las $S_i$). Se define su diferencia $\Phi \equiv \phi_1 - \phi_2$.

Por linealidad, $\Phi$ satisface Laplace en $B$. Y satisface condiciones de borde
\textbf{mucho más simples}: como $\phi_1$ y $\phi_2$ tienen los mismos datos, la
resta los anula. En cada superficie vale una de las dos:

\[
\Phi\big|_{S_i} = 0 \qquad\text{o}\qquad \grad\Phi\cdot\hat n\,\big|_{S_i} = 0
\]

Se define ahora el campo auxiliar $\vec D \equiv \Phi\,\grad\Phi$ y se calcula su
divergencia con la regla del producto
$\div(f\vec A) = \grad f\cdot\vec A + f\,\div\vec A$:

\[
\div\vec D = \grad\Phi\cdot\grad\Phi + \Phi\,\laplacian\Phi = |\grad\Phi|^2
\]

donde el segundo término se anula \textbf{porque $\Phi$ satisface Laplace}.

\begin{notebox}[Interpretación física]
No es necesaria para la prueba, pero sí para entenderla. La densidad de energía
electrostática es $u = \tfrac{1}{2}\varepsilon_0|\vec E|^2$. Como $\grad\Phi$ es
(menos) el campo asociado a $\Phi$, la cantidad $|\grad\Phi|^2$ es proporcional a
la densidad de energía electrostática de la diferencia.
\end{notebox}

Aplicando el teorema de la divergencia sobre $B$, con
$\partial B = S_0 \cup S_1 \cup \cdots \cup S_n$:

\[
\int_B |\grad\Phi|^2\,dV = \int_B \div\vec D\,dV
= \oiint_{\partial B} \Phi\,\grad\Phi\cdot\hat n\,dS = 0
\]

La integral de superficie se anula \textbf{por las condiciones de borde}: en cada
superficie, o bien $\Phi = 0$, o bien $\grad\Phi\cdot\hat n = 0$; en cualquiera de
los dos casos el integrando es nulo ahí.

Queda entonces $\int_B |\grad\Phi|^2\,dV = 0$, con un integrando \textbf{positivo
o nulo y continuo}. Sumando cantidades no negativas y obteniendo cero, todas
tenían que ser cero:

\[
|\grad\Phi|^2 = 0 \;\;\text{en todo }B
\quad\Longrightarrow\quad
\grad\Phi = 0 \;\;\text{en todo }B
\quad\Longrightarrow\quad
\Phi = \text{constante}
\]

\[
\boxed{\;\phi_1 = \phi_2 + \text{cte.}\;}
\]

\begin{warnbox}[La hipótesis de continuidad no es decorativa]
Si la cantidad no fuera suave se podría tener una singularidad en un punto y el
argumento fallaría; se están buscando soluciones suaves.
\end{warnbox}

\begin{keybox}[Lo que está en el corazón de la prueba]
La diferencia entre dos soluciones es una configuración con \textbf{energía
electrostática nula} dentro de $B$. Eso es lo que dice Laplace con condiciones de
borde homogéneas. Nótese que es la \emph{energía de la diferencia}, no la
diferencia de energías; y que $\Phi$ resulta equipotencial aunque el problema
original no tenga por qué serlo.
\end{keybox}

\bigskip
\noindent\emph{Continúa en la Clase 6, con la ecuación de Laplace en otros
sistemas de coordenadas ---cilíndricas, esféricas--- y el comienzo de su
resolución, empezando por las situaciones más simples.}

\end{document}
